% main.tex -- PVLDB Vol. 20 (EA&B track) version of Paper 1, condensed from
% paper/main.tex. Same discipline as Paper 1: every result is a \stat{} macro
% filled from the committed analysis outputs by ../paper/tools/stats_to_tex.py
% (see Makefile); no result is typed. Host-conditional text uses Paper 1's
% "% BEGIN-RAN:<host>" / "% END-RAN:<host>" markers and \if<flag>ran flags.
% Material cut for length is in the extended version (paper/, cited as
% \citet{ejiofor2026extended}); the version lock-in and core-count analyses
% are the companion manuscript's subject (paper2/, \citet{ejiofor2026lockin}).
\documentclass[sigconf, nonacm]{acmart}

%%% do not modify the following VLDB block %%
%%% VLDB block start %%%
\usepackage{pvldb}% VLDB/PVLDB formatting overrides, see pvldb.sty for details.
                      % Do not remove -- required for consistent proceedings
                      % formatting across all VLDB submissions.
%%% VLDB block end %%%

\usepackage{multirow}
\usepackage{enumitem}
\usepackage{seqsplit}  % line-breakable commit hashes and SHA-256 digests
% macros.tex -- statistics lookup, figure inclusion and notice macros.
%
% Numbers in the paper are never typed in. Every result is written as \stat{key},
% where key is a dotted stats.json key (e.g. c1.baseline.sf300.estimate).
% paper/tools/stats_to_tex.py turns <analysis-out>/stats.json into
% generated/stats.tex, which contains one \defstat{key}{value} per key and
% \statsfinaltrue. In final mode a missing key is a hard LaTeX error; in draft
% mode (no generated/stats.tex) the key is printed in a box so the gaps are visible.

\makeatletter
\newif\ifstatsfinal
\statsfinalfalse
\newcommand{\defstat}[2]{\expandafter\def\csname stat@#1\endcsname{#2}}
\newcommand{\statplaceholder}[1]{%
  \fcolorbox{red}{yellow!20}{\scriptsize\ttfamily\detokenize{#1}}}
\newcommand{\stat}[1]{%
  \ifcsname stat@#1\endcsname
    \csname stat@#1\endcsname
  \else
    \ifstatsfinal
      \PackageError{stats}{Missing statistic '#1' in generated/stats.tex}%
        {Re-run tools/stats_to_tex.py; the analysis output lacks this key.}%
    \else
      \statplaceholder{#1}%
    \fi
  \fi}
% "estimate [ci_low, ci_high]" for a ratio prefix such as c1.baseline.sf300
\newcommand{\ratioci}[1]{%
  \stat{#1.estimate}\,{\footnotesize[\stat{#1.ci_low},\,\stat{#1.ci_high}]}}
% "estimate x (95% CI ci_low--ci_high)" for running text
\newcommand{\ratiotext}[1]{%
  \stat{#1.estimate}$\times$ (95\%~CI \stat{#1.ci_low}--\stat{#1.ci_high})}
\makeatother

% GPU-generation ladder (amendments A2, A3). One flag per additional host;
% the H100 is primary and always present. Draft builds show every host's text
% with placeholders; stats_to_tex.py --not-run HOST=REASON writes
% \<flag>ranfalse, and \ladderranfalse when no additional host ran. Wrap
% host text in \if<flag>ran ... \fi inside "% BEGIN-RAN:<host>" /
% "% END-RAN:<host>" comment lines (NOTRUN for the \else case) so the key
% scanner knows which keys each case needs.
\newif\ifaonehundredran  \aonehundredrantrue   % a100
\newif\ifghtwohundredran \ghtwohundredrantrue  % gh200
\newif\ifbtwohundredran  \btwohundredrantrue   % b200
\newif\ifladderran       \ladderrantrue        % any of the three ran
\expandafter\def\csname hostflag@a100\endcsname{aonehundred}
\expandafter\def\csname hostflag@gh200\endcsname{ghtwohundred}
\expandafter\def\csname hostflag@b200\endcsname{btwohundred}

% Disclosure watermarks, set by generated/stats.tex: hand-filled _display values
% (stats_to_tex.py --allow-manual-overrides, or a draft) and synthetic stats.
% Printed diagonally on every page and in the title block.
\newif\ifmanualvalues
\newif\ifsyntheticstats

\IfFileExists{generated/stats.tex}{\input{generated/stats.tex}}{}

\newcommand{\watermarktext}{%
  \ifmanualvalues MANUAL VALUES\fi
  \ifmanualvalues\ifsyntheticstats\ \textperiodcentered\ \fi\fi
  \ifsyntheticstats SYNTHETIC -- NOT RESULTS\fi}
\newcommand{\watermarkpage}{%
  \put(0.5\paperwidth,-0.5\paperheight){\makebox(0,0){\rotatebox{45}{%
    \scalebox{4}{\sffamily\bfseries\textcolor{red!30}{\watermarktext}}}}}}
\AddToHook{shipout/background}{%
  \ifmanualvalues\watermarkpage\else\ifsyntheticstats\watermarkpage\fi\fi}
\newcommand{\watermarknotice}{%
  \ifmanualvalues\\\textcolor{red}{\textbf{\watermarktext}}%
  \else\ifsyntheticstats\\\textcolor{red}{\textbf{\watermarktext}}\fi\fi}

% Host-conditional values for tables: the value if the host's campaign ran,
% \hostnotrun otherwise. Keys: <host>.<key>, <host>.<prefix>.{estimate,...}
% and xhost.<host>.<rest>.{estimate,...}. Not for use inside \if...\fi text
% that may be skipped (they are not conditionals themselves).
\makeatletter
\newcommand{\hostnotrun}{---}
\newcommand{\ifhost@ran}[1]{\csname if\csname hostflag@#1\endcsname ran\endcsname}
\newcommand{\hoststat}[2]{\ifhost@ran{#1}\stat{#1.#2}\else\hostnotrun\fi}
% \hoststator{host}{key}{fallback}: the value if the host ran, else the fallback text.
\newcommand{\hoststator}[3]{\ifhost@ran{#1}\stat{#1.#2}\else#3\fi}
\newcommand{\hostratioci}[2]{\ifhost@ran{#1}\ratioci{#1.#2}\else\hostnotrun\fi}
\newcommand{\xhostratioci}[2]{\ifhost@ran{#1}\ratioci{xhost.#1.#2}\else\hostnotrun\fi}
\makeatother

% Figures are copied into figures/ by tools/sync_figures.py. Until then a
% labelled placeholder box keeps the layout and cross-references working.
\newcommand{\studyfig}[2]{%
  \IfFileExists{figures/#1.pdf}{\includegraphics[width=#2]{figures/#1.pdf}}{%
  \IfFileExists{figures/#1.png}{\includegraphics[width=#2]{figures/#1.png}}{%
  \fbox{\parbox[c][0.28\textheight][c]{0.9\linewidth}{\centering\ttfamily
    figure pending: \detokenize{#1}}}}}}

% Author-supplied detail still missing (affiliation, email, venue, ...).
\newcommand{\todoplaceholder}[1]{\textcolor{red}{[TODO: #1]}}

% Results-not-yet-available marker for draft prose (also flagged with
% "% TODO(results)" comments in the source).
\newcommand{\pending}[1]{\textcolor{gray}{[pending results: #1]}}

% TPC fair-use disclaimer. The template wording is TPC's (Fair Use of TPC
% Benchmarks quick reference, 2019-06-05; current Policies v6.20 sec. 8.1.5).
% The name of the derived work is an open author decision; see paper/README.md.
\newcommand{\workloadname}{H100 study workload}
\newcommand{\tpcdisclaimer}{The \workloadname{} is derived from the TPC-H
  Benchmark and as such is not comparable to published TPC-H Benchmark results,
  as the \workloadname{} results do not comply with the TPC-H Benchmark
  Specification.}
\newcommand{\chartnote}{\emph{Derived from TPC-H; not comparable to published
  TPC-H results (see Section~\ref{sec:scope}).}}

% Engine and contrast names
\newcommand{\eng}[1]{\texttt{#1}}
% Prose names for running text and tables; \id... for the code identifiers
% (arm ids in run configuration and the seeded order assignment).
\newcommand{\polarscpu}{Polars CPU}
\newcommand{\polarsgpu}{Polars GPU}
\newcommand{\sparkcpu}{Spark CPU}
\newcommand{\sparkgpu}{Spark GPU}
\newcommand{\idpolarscpu}{\eng{polars}}
\newcommand{\idpolarsgpu}{\eng{polars\_gpu}}
\newcommand{\idsparkcpu}{\eng{spark\_cpu}}
\newcommand{\idsparkgpu}{\eng{spark\_gpu}}
\newcommand{\q}[1]{\texttt{#1}}

% pvldb-macros.tex -- additions for the PVLDB version only. macros.tex is a
% verbatim copy of paper/macros.tex (the Makefile's "check" target verifies
% this), so the \stat machinery, the host flags and the TPC disclaimer wording
% are exactly Paper 1's.

% Repository of the artifact (public GitHub repository, accepted by the PVLDB
% Vol. 20 submission guidelines as an archival location).
\newcommand{\artifacturl}{https://github.com/ejiod2025/gpu-spark-polars-study}

% Zenodo DOI of the raw run-record archive (dataset record). While empty, no
% DOI sentence is printed (Section "Artifact and extended version").
\newcommand{\zenododoi}{10.5281/zenodo.24125306}
\newcommand{\artifactdoisentence}{%
  \ifdefempty{\zenododoi}{}{ The raw run records of all four hosts are
  archived at \href{https://doi.org/\zenododoi}{doi:\zenododoi}.}}

% Title suffix (Paper 1 rule): only if at least one additional ladder host ran.
\ifladderran
  \newcommand{\laddertitlesuffix}{, Replicated across GPU Generations}
\else
  \newcommand{\laddertitlesuffix}{}
\fi


\renewcommand\vldbdoi{XX.XX/XXX.XX}
\renewcommand\vldbpages{XXX-XXX}
% leave empty if no availability url should be set
\renewcommand\vldbavailabilityurl{\artifacturl}

\setlist{leftmargin=*,itemsep=1pt,topsep=2pt}
\setlength{\emergencystretch}{3em}
% Fixed XMP metadata (acmart loads hyperxmp, which otherwise stamps the wall-clock time and a
% random instance UUID into every build), so that the PDF is byte-reproducible with
% SOURCE_DATE_EPOCH (the Makefile fixes it at 2026-10-30 12:00 UTC).
\hypersetup{pdfdate={2026-10-30T12:00:00Z},
  pdfcreationdate={D:20261030120000Z}, pdfmoddate={D:20261030120000Z},
  pdfmetadate={2026-10-30T12:00:00Z},
  pdfdocumentid={uuid:2e1d0c5a-6a43-4f0e-9c3b-5d1f7b0a9e21},
  pdfinstanceid={uuid:2e1d0c5a-6a43-4f0e-9c3b-5d1f7b0a9e21}}

\begin{document}
\title[CPU Dataframes versus GPU Acceleration on a Single Host]{CPU Dataframes
  versus GPU Acceleration on a Single Host: A Pre-registered, Randomized
  Comparison of Polars and Spark on a TPC-H-Derived
  Workload\laddertitlesuffix{} [Experiment, Analysis \& Benchmark]}

\author{Daniel Ejiofor}
\affiliation{%
  \institution{Boston University}
  \city{Boston}
  \state{MA}
  \country{USA}
}
\email{ejiod@bu.edu}

\begin{abstract}
We compare Polars and Spark, each on the CPU and on the GPU (both GPU paths
use NVIDIA cuDF), on one NVIDIA H100 host. The workload is TPC-H-derived (six
adapted TPC-H queries, SF1--SF300; six complex queries, SF30--SF300), not a
TPC-H benchmark result. Its pre-registered design uses four rounds in a
Williams Latin square, timing-independent integrity gates and 95\% intervals
across rounds. At SF300, GPU speedups over the same framework's CPU engine were
\ratiotext{c1.baseline.sf300} for Polars and \ratiotext{c2.baseline.sf300} for
Spark on the baseline suite, and \ratiotext{c1.complex.sf300} and
\ratiotext{c2.complex.sf300} on the complex suite. The best GPU engine was
faster than Polars CPU from \stat{crossover.baseline.c4.scale} (baseline) and
\stat{crossover.complex.c4.scale} (complex), and slower at
\stat{c4.baseline.summary.below}. Polars GPU rejected
\stat{coverage.polars_gpu.complex.unsupported} of the six complex queries as
unsupported (CPU fallback disabled by design). However, the GPU engine requires
Polars 1.44.x, and Polars 2.0.0 on the CPU was \stat{bridge_summary.both}.
% BEGIN-RAN:any
\ifladderran
Campaigns on \stat{hosts.other_ran_list} showed the same pattern, with GPU
time ratios to the H100 of
% BEGIN-RAN:a100
\ifaonehundredran
\stat{ladder.xhost_gpu.a100.min} to \stat{ladder.xhost_gpu.a100.max} on the
A100 (slower) and
\fi
% END-RAN:a100
\stat{ladder.xhost_gpu.later.min} to \stat{ladder.xhost_gpu.later.max} on the
\stat{ladder.xhost_gpu.later.hosts}.
\fi
% END-RAN:any
Across \stat{ladder.hosts_count} hosts, \stat{ladder.outputs_failed_total} of
\stat{ladder.outputs_compared_total} output comparisons disagreed with Polars CPU.
\end{abstract}

\maketitle

%%% do not modify the following VLDB block %%
%%% VLDB block start %%%
\vldbtopmatter
%%% VLDB block end %%%

\noindent\textbf{Scope notice.} \tpcdisclaimer{} The queries are a selected
and adapted subset with additional custom queries, Float64 normalization of
decimal columns, and our own timing and summary method; see
Section~\ref{sec:scope}.

% ======================================================================
\section{Introduction}
\label{sec:intro}

Practitioners who outgrow a laptop-scale dataframe face two common routes. One
is a faster single-node engine: Polars~\citep{polars1442} executes lazy,
optimized query plans with a multi-threaded and streaming engine written in
Rust. The other is GPU acceleration, either inside the same dataframe API (the
Polars GPU engine built on NVIDIA cuDF~\citep{polars_gpu_guide,cudf_polars_docs})
or underneath an existing Spark SQL workload (NVIDIA's cuDF plugin for Apache
Spark~\citep{nvidia_cudf_spark_docs}). Vendor and project material reports
large speedups for both routes~\citep{semaan2024polarsgpu,nvidia_spark_benchmarks},
usually against a CPU baseline from the same framework.

The decision a practitioner actually faces is broader: \emph{for this kind of
query at this data size, is a GPU engine faster than the best CPU engine I could
use on the same machine?} Answering that requires holding the host, the input
files, the query semantics and the timing boundary fixed across all engines,
checking that every engine returns the same answer, and checking that a ``GPU''
timing really ran on the GPU. It also requires enough repeated, order-balanced
measurement to separate engine effects from run-to-run noise and order
effects~\citep{mytkowicz2009wrong,kalibera2013rigorous,hoefler2015scientific,raasveldt2018fair}.

An earlier study by the author compared three of these engines (Polars CPU,
Spark CPU, and Spark with the GPU plugin) on an A100 40\,GB host. It found that
the plugin accelerated Spark on the larger scales but that Polars CPU retained
the lower aggregate score at every observed scale. That study ran one session
per engine and scale, so it had no replication and no interval estimates, and
it did not include the Polars GPU engine. This paper replaces it with a
pre-registered, replicated design.

\paragraph{Contributions.}
\begin{enumerate}
  \item A four-engine comparison (CPU and GPU variants of Polars and Spark) on
    one H100 host with identical, verified inputs, from SF1 to SF300
    (1.8~billion \texttt{lineitem} rows).
  \item A pre-registered design (Section~\ref{sec:design}): a Williams Latin
    square over four rounds, fixed contrasts, and a statistical plan that
    reports geometric-mean ratios with across-round confidence intervals rather
    than single-run point values.
  \item Integrity gates that are independent of timings: cross-engine output
    agreement, a strict no-CPU-fallback guard for the Polars GPU engine, and
    final-plan checks for the Spark plugin.
  \item A GPU coverage analysis that reports which queries cannot run on each
    GPU engine, rather than silently timing a CPU fallback.
  \item A side study measuring how much the Polars version (1.44.2 versus
    2.0.0) changes the CPU baseline that the GPU speedup is measured against.
% BEGIN-RAN:any
\ifladderran
  \item A GPU-generation ladder: the complete campaign repeated, with the same
    design, software versions, inputs and engine thread count, on
    \stat{hosts.other_ran_list}, each analyzed separately and compared
    descriptively with the H100 (Section~\ref{sec:ladder}).
\fi
% END-RAN:any
  \item A complete artifact: code, saved outputs, plans, telemetry and
    manifests, from which every number in this paper is regenerated
    (Section~\ref{sec:artifact}).
\end{enumerate}

\paragraph{Summary of findings.} On the H100, both GPU engines lost to their
own framework's CPU engine at the smallest scale (C1 below~1 at
\stat{c1.baseline.summary.below}, C2 at \stat{c2.baseline.summary.below}), and
the best GPU engine beat Polars CPU from \stat{crossover.baseline.c4.scale}
(baseline suite), by multiples of a few rather than orders of magnitude. The
Polars GPU engine was the faster GPU engine in \stat{ladder.c3.cells_above} of
\stat{ladder.c3.cells} suite--scale cells across all hosts, but could not run
\stat{coverage.polars_gpu.complex.unsupported} of the twelve queries; the
Spark plugin ran all of them. The CPU baseline that
makes these speedups possible is version-locked: the GPU engine needs an older
Polars, and the current Polars release was substantially faster on the CPU.
% BEGIN-RAN:any
\ifladderran
Repeating the campaign on older and newer GPUs did not change the decision
pattern, and the GPUs newer than the H100 did not make the GPU engines
consistently faster.
\fi
% END-RAN:any
This paper is a condensed version of a longer report~\citep{ejiofor2026extended},
which gives every design detail, every secondary analysis and per-host
appendices; the version lock-in and the core-count confounds of the ladder are
analyzed in depth in a companion manuscript~\citep{ejiofor2026lockin}.

% ======================================================================
\section{Background and Related Work}
\label{sec:background}

\paragraph{Dataframe engines.} The pandas library~\citep{mckinney2010pandas}
made the dataframe the default interface for tabular analysis in Python;
\citet{petersohn2020dataframes} discuss why its eager, single-threaded execution
model is hard to scale. Polars~\citep{polars1442,polars200} provides a lazy API
whose query plans are optimized (predicate and projection pushdown, among other
rewrites) and executed by a multi-threaded engine; its streaming engine
processes data in batches. DuckDB~\citep{raasveldt2019duckdb} is an embedded
analytical SQL engine; here it is used only to generate data.

\paragraph{GPU dataframes and Spark.} cuDF is NVIDIA's GPU dataframe library,
built on the C++ library libcudf. The Polars GPU engine (\texttt{cudf-polars})
takes an optimized Polars plan and executes it with cuDF; when a plan contains
an unsupported operation it falls back to the CPU engine unless the user sets
\texttt{raise\_on\_fail=True}~\citep{polars_gpu_guide,cudf_polars_docs}. The
Polars user guide states that the engine is in open beta and that speedups are
most likely for workloads dominated by grouped aggregations and joins, and that
I/O-bound queries perform about the same as on the CPU~\citep{polars_gpu_guide}.
Each \texttt{cudf-polars} release supports a bounded range of Polars versions.
Apache Spark generalizes MapReduce with resilient distributed
datasets~\citep{zaharia2012rdd,zaharia2016spark}; Spark SQL adds the Catalyst
optimizer and the DataFrame API~\citep{armbrust2015sparksql}, and adaptive query
execution re-optimizes plans at runtime~\citep{spark359docs}. NVIDIA's cuDF
plugin for Spark~\citep{nvidia_cudf_spark_docs,nvidia_cudf_spark_repo} replaces
supported physical operators with GPU implementations and leaves the rest on
the CPU, so a GPU-enabled Spark plan can be hybrid. We therefore inspect final
plans rather than assume GPU execution.

\paragraph{GPU and CPU analytics.} \citet{bress2014gpudb} and
\citet{rosenfeld2022heterogeneous} survey GPU-accelerated and heterogeneous
CPU/GPU query processing. \citet{shanbhag2020gpucpu} argue that because
analytical operators are usually memory-bandwidth bound, GPU speedups over
well-implemented CPU code should be expected near the GPU-to-CPU bandwidth
ratio rather than orders of magnitude. \citet{lutz2020pump} show that the
CPU--GPU transfer link bottlenecks GPU query processing on data larger than
device memory, and that faster interconnects relieve it. Public GPU dataframe
comparisons are mostly vendor or project benchmarks: the Polars GPU
announcement reports speedups over Polars CPU on PDS-H using an
H100~\citep{semaan2024polarsgpu}, and the database-like ops
benchmark~\citep{dbbenchmark} compares CPU tools on grouping and joins. To our
knowledge, independent, replicated comparisons of a GPU dataframe engine and a
GPU Spark plugin against a strong single-node CPU engine on identical inputs
remain scarce.

\paragraph{TPC-H and derived use.} TPC-H is the TPC's decision-support
benchmark~\citep{tpch301,poess2000tpc}; \citet{boncz2014tpch} analyze the
``choke points'' its queries exercise. Many evaluations reuse its schema,
generator and queries without following the full specification, under names
that mark the derivation, such as the Polars project's PDS-H~\citep{vink2025pdsh}
and NVIDIA's NDS-H~\citep{nvidia_spark_benchmarks}. We generate data with
DuckDB's \texttt{tpch} extension, which documents small differences from the
official specification~\citep{duckdb_tpch_ext}. TPC's fair-use material
requires such studies to carry a prescribed disclaimer, disclose deviations and
the test environment, and avoid comparison with official
results~\citep{tpcfairuse2019,tpcpolicies620}.

\paragraph{Benchmarking methodology.} Hidden factors such as environment
layout can bias measurements~\citep{mytkowicz2009wrong}; repeated, randomized
execution with interval estimates is recommended~\citep{kalibera2013rigorous,
hoefler2015scientific}, and database comparisons have many fairness
pitfalls~\citep{raasveldt2018fair}. Ratios of times should be summarized with
the geometric mean~\citep{fleming1986geomean}. Williams designs balance
first-order carryover in crossover experiments~\citep{williams1949latin}.

% ======================================================================
\section{Experimental Design}
\label{sec:design}

\subsection{Pre-registration and amendments}
\label{sec:prereg}

The research questions, arms, workload, order design, integrity gates,
pre-declared GPU gaps and statistical plan below were locked on 2026-10-29,
before any timed data existed on any host, in the design document of the study
repository, \texttt{docs/H100\_EXPERIMENT\_DESIGN.md}. Locked sections are never edited; changes are
appended as dated amendments with their reason. A1 separates \emph{coverage
outcomes} (reported) from \emph{integrity failures} (which fail a round;
Section~\ref{sec:gates}). A2 and A3 add the B200, A100 and GH200 hosts, each
running its own complete campaign under the unchanged design
(Section~\ref{sec:hosts}). A4 records a known CUDA 12.9 compiler defect on
Blackwell that can silently corrupt \texttt{thrust::exclusive\_scan\_by\_key}
results~\citep{cccl11167}, for which the Polars GPU build carries cuDF's
workaround~\citep{cudf23973}; it changes no gate or analysis. A1--A3 preceded
all timed data and A4 all B200 data. No decision about the design, stopping or
analysis was based on observed speeds.

The locked design with A1--A3 was committed as \texttt{ee2e6a7} (full hash
\texttt{\seqsplit{ee2e6a76a0f0b8ddd28bc4a9293dde5176c34ffe}}) and pushed on
2026-10-30 (commit time 03:55\,UTC); the design file at that commit has SHA-256
\texttt{\seqsplit{6b7544ee23b4ad40d66c8ba0866772ce0fec0ca6a6db92b6435f064979dcddff}}.
A4 was committed separately as \texttt{afd07c0}. Each campaign records the
SHA-256 of the design document it ran under; for the H100 this resolves to
commit \texttt{\stat{meta.design_commit_short}} with \stat{meta.amendments}
amendments (added later: \stat{meta.later_amendments_list}), and the campaign
started at \stat{meta.campaign_start_utc}, after that commit. Every other
campaign also started after the commit of the design version it records.

\subsection{Research questions}
\label{sec:rqs}
\textbf{RQ1}: How much does the GPU accelerate each framework (Polars GPU
versus Polars CPU; Spark GPU versus Spark CPU)? \textbf{RQ2}: Which engine is
fastest per query, suite and scale, and does any GPU engine beat Polars CPU?
\textbf{RQ3}: At which scale from SF1 to SF300, if any, does a GPU engine
overtake the best CPU engine? \textbf{RQ4}: Which queries cannot run on each
GPU engine, and why? \textbf{RQ5} (side study): How much does the Polars
version (1.44.2 versus 2.0.0, CPU) change Polars CPU timings?

\subsection{Arms}
\label{sec:arms}

\begin{table}[t]
\centering
\caption{The four arms. No arm was tuned during the campaign; the Spark
configurations are byte-identical to the earlier A100 study.}
\label{tab:arms}
\footnotesize
\setlength{\tabcolsep}{3pt}
\begin{tabular}{@{}lp{0.70\linewidth}@{}}
\toprule
Arm & Engine, version and settings \\
\midrule
\polarscpu{} & Polars 1.44.2 streaming CPU engine, Python 3.11, 26 threads \\
\polarsgpu{} & Polars 1.44.2 + \texttt{cudf-polars-cu12} 26.10.2; \texttt{raise\_on\_fail=True}; streaming executor (defaults) \\
\sparkcpu{}  & PySpark 3.5.9, \texttt{local[26]}, Java 17, 96\,GiB driver heap, 120 shuffle partitions, AQE on \\
\sparkgpu{}  & As \sparkcpu{} plus \texttt{rapids-4-spark\_2.12} 26.08.1 (SHA-256 recorded), 2 concurrent GPU tasks, 2\,GB pinned pool \\
\bottomrule
\end{tabular}
\end{table}

Table~\ref{tab:arms} lists the arms. Polars runs at version 1.44.2 rather than
the current 2.0.0 because no released \texttt{cudf-polars} supports Polars 2.0:
the newest release, 26.10.2, requires Polars ${\ge}1.35,{<}1.45$. Both Polars
arms share the same Polars build, so the RQ1 contrast is not confounded by
version; RQ5 measures the version effect on the CPU separately. Both Spark arms
also set 128\,MiB maximum partition bytes and the UTC session time zone, and
the GPU arm enables the plugin's variable floating-point aggregation and logs
operators not placed on the GPU. The Spark settings are the earlier study's
fixed values, not tuned ones; the Polars GPU
engine uses \texttt{cudf-polars}' default streaming-executor options and memory
resource, recorded per session. Thread counts set engine worker pools
(\texttt{POLARS\_MAX\_THREADS=26}, \texttt{local[26]}).

\subsection{Hosts and the GPU-generation ladder}
\label{sec:hosts}

\begin{table*}[t]
\centering
\caption{Hosts (Lambda on-demand instances), as recorded by each campaign. The
H100 is the primary analysis; each other host runs its own complete campaign
and is never pooled with it. vCPUs: the count the host reports, then the 26
engine threads and the physical cores they occupy (\texttt{lscpu} topology,
inside the virtual machine). Driver: version and, in parentheses, the highest
CUDA version it supports; every host runs the CUDA \stat{hosts.cuda_runtime}
runtime. \hostnotrun{}: campaign not run.}
\label{tab:hosts}
\small
\setlength{\tabcolsep}{4pt}
\begin{tabular}{@{}llrlllll@{}}
\toprule
Host & Instance type & \shortstack[r]{GPU\\memory} & ISA & \shortstack[l]{vCPUs (threads\\on cores)} & RAM & Driver (CUDA) & CPU model \\
\midrule
A100  & \texttt{gpu\_1x\_a100\_sxm4} & \hoststat{a100}{host.gpu_memory_gib}\,GiB & x86\_64 & \hoststat{a100}{host.cpu_layout} & \hoststat{a100}{host.host_ram_gib}\,GiB & \hoststat{a100}{meta.driver_cuda} & \hoststat{a100}{meta.cpu_model} \\
H100  & \texttt{gpu\_1x\_h100\_sxm5} & \stat{host.gpu_memory_gib}\,GiB & x86\_64 & \stat{host.cpu_layout} & \stat{host.host_ram_gib}\,GiB & \stat{meta.driver_cuda} & \stat{meta.cpu_model} \\
GH200 & \texttt{gpu\_1x\_gh200} & \hoststat{gh200}{host.gpu_memory_gib}\,GiB & aarch64 & \hoststat{gh200}{host.cpu_layout} & \hoststat{gh200}{host.host_ram_gib}\,GiB & \hoststat{gh200}{meta.driver_cuda} & \hoststat{gh200}{meta.cpu_model} \\
B200  & \texttt{gpu\_1x\_b200\_sxm6} & \hoststat{b200}{host.gpu_memory_gib}\,GiB & x86\_64 & \hoststat{b200}{host.cpu_layout} & \hoststat{b200}{host.host_ram_gib}\,GiB & \hoststat{b200}{meta.driver_cuda} & \hoststat{b200}{meta.cpu_model} \\
\bottomrule
\end{tabular}
\end{table*}

All engines ran sequentially on one Lambda \texttt{gpu\_1x\_h100\_sxm5}
instance: one NVIDIA H100 80\,GB SXM5 (\stat{meta.gpu_name}, driver
\stat{meta.driver_version}), \stat{host.vcpus} vCPUs (\stat{meta.cpu_model}),
\stat{host.host_ram_gib}\,GiB RAM and a local NVMe volume holding the data and
Spark's local directory. GPU clocks were not locked; they were recorded.

Three further hosts (Table~\ref{tab:hosts}) place the H100 on a ladder of GPU
generations: an A100 40\,GB (Ampere; the GPU of the earlier study), a GH200
(Grace~Hopper: a 96\,GB Hopper GPU joined to an Arm-based Grace CPU by a
coherent NVLink-C2C link~\citep{nvidia_gh200}) and a B200 (Blackwell). The
provider offered no newer GPU. Each host runs its own complete campaign under
the unchanged design (same seed, Latin square, gates and side study), with its
inputs regenerated by the same recipe and passing the same input gate; data are
never pooled, and every primary estimate is computed per host. All hosts run
the engines with 26 threads, so engine thread counts are equal across the
ladder and further vCPUs are left idle. Equal threads are not equal physical
cores: on the H100, with \stat{host.smt_threads_per_core} hardware threads per
core, the 26 threads occupy \stat{host.cpus_arg_cores_text} physical cores.
Every host runs the pinned CUDA runtime with its provider image's driver,
through CUDA minor-version
compatibility~\citep{nvidia_cuda_compat} where the driver is older (here:
\stat{hosts.minor_compat_list}). A host's campaign runs only if the provider
has capacity and the pinned stack passes preflight and input verification on
it; versions are never changed to make a host run. On the GH200 (aarch64) the
same package versions are used, with the platform-specific artifacts recorded
by SHA-256.
% BEGIN-NOTRUN:a100
\ifaonehundredran\else
The A100 campaign was not run: \stat{a100.not_run_reason}.
\fi
% END-NOTRUN:a100
% BEGIN-NOTRUN:gh200
\ifghtwohundredran\else
The GH200 campaign was not run: \stat{gh200.not_run_reason}.
\fi
% END-NOTRUN:gh200
% BEGIN-NOTRUN:b200
\ifbtwohundredran\else
The B200 campaign was not run: \stat{b200.not_run_reason}.
\fi
% END-NOTRUN:b200

\subsection{Workload and inputs}
\label{sec:workload}

\begin{table}[t]
\centering
\caption{Queries. Baseline queries are adapted from DuckDB 1.5.6's bundled
TPC-H query text; \q{cx\_q09} and \q{cx\_q13} are derived from TPC-H Q9 and
Q13; the other complex queries are custom. Spark runs the SQL; Polars runs
equivalent native lazy expressions.}
\label{tab:queries}
\footnotesize
\setlength{\tabcolsep}{3pt}
\begin{tabular}{@{}llp{0.62\linewidth}@{}}
\toprule
Suite & Query & Main work \\
\midrule
\multirow{6}{*}{baseline}
 & \q{q01} & Broad ship-date filter; grouped sums, averages, counts \\
 & \q{q03} & 3-table join, selective predicates, revenue per order, top 10 \\
 & \q{q06} & Selective filter, single revenue sum \\
 & \q{q10} & 4-table join, customer grouping with strings, top 20 \\
 & \q{q12} & 2-table join, conditional counts by ship mode \\
 & \q{q14} & 2-table join, promotional revenue share \\
\midrule
\multirow{6}{*}{complex}
 & \q{cx\_q09} & 6-table join incl.\ composite key, substring filter \\
 & \q{cx\_q13} & Left outer join, comment pattern filter, nested aggregation \\
 & \q{cx\_semi\_anti} & Correlated \texttt{EXISTS}/\texttt{NOT EXISTS} (semi/anti joins) \\
 & \q{cx\_window} & \texttt{LAG} and \texttt{ROW\_NUMBER} over two orderings \\
 & \q{cx\_distinct} & 4-table join, three exact distinct counts \\
 & \q{cx\_self\_join} & \texttt{lineitem} self-join with residual inequality \\
\bottomrule
\end{tabular}
\end{table}

Table~\ref{tab:queries} lists the twelve queries. The baseline suite runs at
SF1, 10, 30, 100 and 300 with the earlier study's session lifecycle. The
complex suite runs at SF30, 100 and 300 with a \texttt{verified\_gc} barrier
after every query, outside the timer: garbage collection (Python, and the JVM
for Spark), then waiting until the Spark session owns no shuffle files and has
no active jobs; a failed barrier aborts the session. Extra sort keys make
top-$N$ ties deterministic. Decimal quantity and money columns are cast to
Float64 by every engine, so exact decimal arithmetic is not measured.

Inputs are generated on each host with DuckDB 1.5.6's \texttt{tpch}
extension~\citep{duckdb_tpch_ext} with the earlier study's parameters
(partitioned generation for SF100 and SF300) and written as Snappy Parquet; SF300
has 1,799,989,091 \texttt{lineitem} rows. Before round~1 a verifier must
report that row counts, schemas, generation parameters and order-independent
content fingerprints match the earlier study's manifests (H100 gate result:
\stat{meta.input_gate}). Byte identity is not required: DuckDB's per-thread
output makes the assignment of rows to files depend on scheduling, so file
hashes differ between regenerations while table contents do not.

\subsection{Sessions, rounds and order}
\label{sec:structure}

A \emph{session} is one engine process for one (suite, scale). It preheats the
input files (untimed), runs one untimed warmup per query, and then $R=3$
measured trials per query in a seeded, shuffled order. The timer covers query
construction, optimization, Parquet reading, execution and result collection,
including the device-to-host transfer for GPU engines; for Polars it stops when
\texttt{collect} returns a host DataFrame, for Spark when \texttt{collect()}
returns Python \texttt{Row} objects. A \emph{round} runs every engine on every
(suite, scale) once (32 sessions); the \emph{campaign} is four rounds, and the
round is the block for analysis.

Engine order follows a Williams balanced Latin square~\citep{williams1949latin}
for four engines, with letters assigned by a seeded shuffle (seed 20261029):
A~=~\sparkcpu{}, B~=~\polarsgpu{}, C~=~\polarscpu{}, D~=~\sparkgpu{}, and
rounds ABDC, BCAD, CDBA and DACB. Each engine takes each position once and
each ordered pair of engines is adjacent exactly once, so first-order carryover
(page-cache or thermal state left by the previous engine) is balanced. Suite
order alternates between rounds and scales ascend within a suite. No other
workload ran on the host; before each session no GPU process was running and
the previous session's processes had exited. GPU utilization, memory, power,
temperature and clocks were sampled every 200\,ms and host resources every
second.

\subsection{Integrity gates and coverage}
\label{sec:gates}

Each round must pass all of the following, or it is repeated once (and the
repeat reported); a second failure stops the campaign. None depends on
timings. (1)~Every session exits, with no timeout and no unclassified error.
(2)~Every output agrees with \polarscpu{} for the same query (relative
tolerance $10^{-8}$, absolute $10^{-6}$, integers exact). (3)~For
\polarsgpu{}, a strict session guard confirms the resolved executor and
\texttt{raise\_on\_fail=True} and runs a canary query with a Python
user-defined function, which must raise; otherwise the engine could be falling
back to the CPU silently. (4)~For \sparkgpu{}, final plans contain GPU scan and
compute operators and no CPU core operators. (5)~Every \texttt{verified\_gc}
barrier passes.

Under amendment A1 failures are classified, and only integrity failures fail a
round. A \emph{coverage outcome} is a GPU engine rejecting or failing to
complete a query because of (a) an unsupported operation (for example a
\texttt{cudf-polars} \texttt{NotImplementedError} under
\texttt{raise\_on\_fail=True}) or (b) resource exhaustion (device or host out of
memory). Coverage outcomes are reported under RQ4 and excluded from contrasts
involving that engine; a coverage outcome in some rounds but not others makes
the cell \emph{unstable} and excluded in all rounds. An \emph{integrity
failure} is an output disagreement, a silent CPU fallback, CPU core operators
in a GPU plan, a failed barrier, a timeout, or any error matching no fixed
classification rule. Preflight on 2026-10-26 showed that \polarsgpu{} rejects
\q{cx\_q13} (the inline regular-expression flag \texttt{(?s)}) and
\q{cx\_window} (\texttt{cum\_count().over(...)}); these were pre-declared as
coverage outcomes, and queries are never rewritten for the GPU.

\subsection{Statistical analysis plan}
\label{sec:stats}

Let $t_{e,q,r}$ be the median of the three measured trials of engine $e$ on
query $q$ in round $r$. Each contrast compares a reference engine $a$ with an
engine $b$ as the speedup ratio $\rho_{q,r} = t_{a,q,r} / t_{b,q,r}$, so
$\rho > 1$ means $b$ is faster: \textbf{C1}~\polarscpu{}\,/\,\polarsgpu{} and
\textbf{C2}~\sparkcpu{}\,/\,\sparkgpu{} (RQ1); \textbf{C3}~\sparkgpu{}\,/\,\polarsgpu{}
(GPU head-to-head, RQ2); and \textbf{C4}~\polarscpu{}\,/\,best GPU engine, the
faster of the GPU engines that ran the query in that round (RQ2, RQ3). With
$\ell_{q,r} = \log \rho_{q,r}$, the per-query estimate is $\exp(\bar\ell_q)$,
the geometric mean over the four rounds, with 95\% interval
$\exp\!\bigl(\bar\ell_q \pm t_{0.975,3}\, s_q / \sqrt{4}\bigr)$,
$t_{0.975,3} \approx 3.182$. Per (suite, scale), the per-round suite statistic
is the mean of $\ell_{q,r}$ over the queries supported by both engines, with the
same $t$-interval across rounds (the round is the block). The \textbf{RQ3
crossover} is the smallest scale, per suite, at which the C4 interval lies
entirely above~1. Contrasts are reported as estimates with intervals; no
hypothesis tests or $p$-values are used for claims, and with $\text{df}=3$
intervals are expected to be wide.

\textbf{Side study (RQ5).} After the main rounds, on the same host, \polarscpu{}
1.44.2 is compared with Polars 2.0.0 (the earlier study's version) in four
rounds ordered AB, BA, BA, AB, with the same session protocol; the estimate is
$t_{2.0.0} / t_{1.44.2}$ per suite and scale. \textbf{Cross-host analysis}
(secondary, descriptive): for each host $h$, engine, suite and scale, the
geometric mean over queries of $t_{h}/t_{\mathrm{H100}}$ (below~1: faster on
$h$), with a Welch-type interval on the log scale. For GPU engines this
confounds the GPU with the CPU, memory, CPU--GPU link and driver.

\subsection{Scope, deviations and notice}
\label{sec:scope}
\tpcdisclaimer{} Deviations from the TPC-H specification include: a selected
subset of six adapted queries with fixed substitution parameters, plus two
TPC-H-derived and four custom complex queries; data from DuckDB's generator,
which differs slightly from the specification~\citep{duckdb_tpch_ext};
Float64 instead of decimal arithmetic; extra sort keys; Parquet files rather
than a loaded database; no refresh functions, power or throughput tests, or
TPC metrics; and our own timing boundary and summary statistics. No result in
this paper may be compared with official TPC-H results. Use of the benchmark
name and the TPC software is subject to the TPC
policies~\citep{tpcpolicies620,tpcfairuse2019} and the TPC
EULA~\citep{tpceula22}.
% TODO(author): confirm TPC written-permission status for benchmark-name use
% (Policies 6.20 sec. 8.1.4) and the final name of the derived workload
% (\workloadname in macros.tex). See paper/README.md.

% ======================================================================
\section{Results on the H100}
\label{sec:results}

\paragraph{Campaign integrity.} \stat{meta.rounds_completed} rounds completed
(\stat{meta.rounds_excluded_count} excluded, \stat{meta.rounds_repeated}
repeated under the gate rule) with \stat{meta.sessions_completed} sessions and
\stat{meta.measured_trials} measured trials. All \stat{meta.outputs_compared}
saved outputs were compared with \polarscpu{}; \stat{meta.outputs_failed}
failed (largest relative difference \stat{meta.max_rel_diff}, against a
tolerance of $10^{-8}$). \stat{meta.spark_gpu_plans_checked} \sparkgpu{} final
plans passed the plan rule, and \stat{meta.cleanup_barriers_passed}
\texttt{verified\_gc} barriers passed. Campaign window:
\stat{meta.campaign_start_utc} to \stat{meta.campaign_end_utc}.

\begin{table*}[t]
\centering
\caption{H100: suite-level geometric-mean speedup ratios with 95\% CIs ($t$,
df~=~3, across four rounds); $>1$: the second-named engine is faster. C1 and
C3 exclude \q{cx\_q13} and \q{cx\_window} in the complex suite. Right column:
the Polars version bridge (RQ5), $t_{2.0.0}/t_{1.44.2}$ on the CPU; below~1:
Polars~2.0.0 is faster. \chartnote}
\label{tab:contrasts}
\small
\setlength{\tabcolsep}{4.5pt}
\begin{tabular}{@{}llccccc@{}}
\toprule
Suite & Scale & C1 \polarscpu{}/\polarsgpu{} & C2 \sparkcpu{}/\sparkgpu{} & C3 \sparkgpu{}/\polarsgpu{} & C4 \polarscpu{}/best GPU & Bridge $t_{2.0.0}/t_{1.44.2}$ \\
\midrule
baseline & SF1   & \ratioci{c1.baseline.sf1}   & \ratioci{c2.baseline.sf1}   & \ratioci{c3.baseline.sf1}   & \ratioci{c4.baseline.sf1}   & \ratioci{bridge.baseline.sf1} \\
         & SF10  & \ratioci{c1.baseline.sf10}  & \ratioci{c2.baseline.sf10}  & \ratioci{c3.baseline.sf10}  & \ratioci{c4.baseline.sf10}  & \ratioci{bridge.baseline.sf10} \\
         & SF30  & \ratioci{c1.baseline.sf30}  & \ratioci{c2.baseline.sf30}  & \ratioci{c3.baseline.sf30}  & \ratioci{c4.baseline.sf30}  & \ratioci{bridge.baseline.sf30} \\
         & SF100 & \ratioci{c1.baseline.sf100} & \ratioci{c2.baseline.sf100} & \ratioci{c3.baseline.sf100} & \ratioci{c4.baseline.sf100} & \ratioci{bridge.baseline.sf100} \\
         & SF300 & \ratioci{c1.baseline.sf300} & \ratioci{c2.baseline.sf300} & \ratioci{c3.baseline.sf300} & \ratioci{c4.baseline.sf300} & \ratioci{bridge.baseline.sf300} \\
\midrule
complex  & SF30  & \ratioci{c1.complex.sf30}   & \ratioci{c2.complex.sf30}   & \ratioci{c3.complex.sf30}   & \ratioci{c4.complex.sf30}   & \ratioci{bridge.complex.sf30} \\
         & SF100 & \ratioci{c1.complex.sf100}  & \ratioci{c2.complex.sf100}  & \ratioci{c3.complex.sf100}  & \ratioci{c4.complex.sf100}  & \ratioci{bridge.complex.sf100} \\
         & SF300 & \ratioci{c1.complex.sf300}  & \ratioci{c2.complex.sf300}  & \ratioci{c3.complex.sf300}  & \ratioci{c4.complex.sf300}  & \ratioci{bridge.complex.sf300} \\
\bottomrule
\end{tabular}
\end{table*}

\begin{figure*}[t]
\centering
\studyfig{f10_headline_times}{0.82\textwidth}
\caption{H100: per-query median time for each engine at SF300 (geometric mean
over four rounds of per-round medians). Bars span the fastest to slowest round.
A query with no timing on an engine appears as a hollow marker in the hatched
``no timing'' strip, labelled with the reason (here: not supported). Log scale;
lower is better. \chartnote}
\Description{Two panels, baseline and complex suite at SF300, with one marker
per engine and query on a logarithmic time axis.}
\label{fig:f10}
\end{figure*}

\subsection{RQ1: GPU speedup within each framework}
\label{sec:rq1}

\begin{figure*}[t]
\centering
\studyfig{f01_forest_sf300}{0.80\textwidth}
\caption{H100: per-query speedup ratios at SF300 (geometric mean over four
rounds) with 95\% CIs, by contrast (C1--C4) and suite; the other scales are in
the extended version. Vertical line at~1: no difference; right of~1: the
second-named engine is faster. Log scale. Intervals have df~=~3 and are wide by
design. \chartnote}
\Description{Forest plots of per-query speedup ratios with confidence
intervals for the four contrasts, baseline and complex suites at SF300.}
\label{fig:f1}
\end{figure*}

Table~\ref{tab:contrasts} gives the suite-level contrasts. On the baseline
suite at SF300, the Polars GPU engine's geometric-mean speedup over Polars CPU
was \ratiotext{c1.baseline.sf300}; the Spark plugin's speedup over CPU Spark was
\ratiotext{c2.baseline.sf300}. On the complex suite at SF300 the corresponding
values were \ratiotext{c1.complex.sf300} (over
\stat{c1.complex.sf300.n_queries} supported queries) and
\ratiotext{c2.complex.sf300}. At SF1, C1 was \ratiotext{c1.baseline.sf1} and C2
was \ratiotext{c2.baseline.sf1}.

On the baseline suite, C1 \stat{c1.baseline.summary.trend}; its interval lay
below~1 at \stat{c1.baseline.summary.below} (Polars CPU faster) and above~1 at
\stat{c1.baseline.summary.above}. On the complex suite C1 ranged from
\stat{c1.complex.summary.min} to \stat{c1.complex.summary.max} and its interval
lay above~1 at \stat{c1.complex.summary.above}. C2 behaved similarly: on the
baseline suite its interval was \stat{c2.baseline.summary.pattern}; on the
complex suite it ranged from \stat{c2.complex.summary.min} to
\stat{c2.complex.summary.max}. Comparing suites,
\stat{c12.complex_larger_text}. Per query (Figure~\ref{fig:f1}), the smallest C1 at SF300 was the
selective scan-and-sum \q{q06} (\ratiotext{c1.baseline.sf300.q.q06}; interval
\stat{c1.baseline.sf300.q.q06.verdict}), against
\ratiotext{c1.baseline.sf300.q.q01} for the grouped aggregation \q{q01} and
\ratiotext{c1.baseline.sf300.q.q03} for the three-table join \q{q03}.

\subsection{RQ2: Fastest engine per query, suite and scale}
\label{sec:rq2}

Across the \stat{fastest.total.baseline.count} baseline query--scale cells, the
fastest engine (lowest geometric-mean median over rounds) was \polarscpu{} in
\stat{fastest.polars.baseline.count}, \polarsgpu{} in
\stat{fastest.polars_gpu.baseline.count}, \sparkcpu{} in
\stat{fastest.spark_cpu.baseline.count} and \sparkgpu{} in
\stat{fastest.spark_gpu.baseline.count}. Across the
\stat{fastest.total.complex.count} complex cells the counts were
\stat{fastest.polars.complex.count}, \stat{fastest.polars_gpu.complex.count},
\stat{fastest.spark_cpu.complex.count} and
\stat{fastest.spark_gpu.complex.count}. \polarscpu{} won its baseline cells at
\stat{fastest.polars.baseline.won_text}. \polarsgpu{} won
\stat{fastest.polars_gpu.complex.of_supported} complex cells it could run; the
complex cells it did not win were won by \sparkgpu{}
(\stat{fastest.spark_gpu.complex.won_text}) and \polarscpu{}
(\stat{fastest.polars.complex.won_text}). \sparkcpu{} won no cell. A
``fastest'' count does not imply a resolved difference.
Figure~\ref{fig:f10} shows the SF300 times behind these counts.

Head-to-head, the two GPU engines' C3 was \ratiotext{c3.baseline.sf300} on the
baseline suite and \ratiotext{c3.complex.sf300} on the complex suite at SF300;
it ranged from \stat{c3.baseline.summary.min} to \stat{c3.baseline.summary.max}
on the baseline suite (it \stat{c3.baseline.summary.trend}) and from
\stat{c3.complex.summary.min} to \stat{c3.complex.summary.max} on the complex
suite, with intervals above~1 \stat{c3.both.above_text}. Against Polars CPU,
the best GPU engine per query (C4) reached \ratiotext{c4.baseline.sf300} and
\ratiotext{c4.complex.sf300} at SF300.

\subsection{RQ3: Crossover scale}
\label{sec:rq3}

\begin{figure*}[t]
\centering
\studyfig{f02_suite_speedup_vs_scale}{0.86\textwidth}
\caption{H100: suite geometric-mean speedup versus scale factor for C1--C4,
with 95\% CI bands (df~=~3). Lines connect measured scales only; they are not a
fitted model and do not extrapolate. \chartnote}
\Description{Line charts of suite speedup against scale factor for the four
contrasts, with confidence bands and a reference line at 1.}
\label{fig:f2}
\end{figure*}

By the pre-registered rule, the baseline crossover scale was
\stat{crossover.baseline.c4.scale} and the complex crossover scale was
\stat{crossover.complex.c4.scale} (Figure~\ref{fig:f2}). Below the crossover,
C4's interval lay below~1 at \stat{c4.baseline.summary.below} on the baseline
suite, so there Polars CPU was faster than either GPU engine. On the complex
suite the interval lay above~1 at \stat{c4.complex.summary.above}. A crossover
at a suite's smallest measured scale (SF30 for the complex suite) only says
that the GPU was already ahead there, and a crossover between two measured
scales is not localized more precisely than that interval.

\subsection{RQ4: GPU coverage}
\label{sec:rq4}

\begin{figure*}[t]
\centering
\studyfig{f04_gpu_coverage}{0.62\textwidth}
\caption{H100: GPU coverage by query, scale and GPU engine (amendment A1
classes): ran on the GPU with all gates passed; not supported; resource
exhaustion; other error (an integrity failure); or partial (unstable coverage).}
\Description{A matrix of queries by scale for each GPU engine, coloured by
coverage outcome.}
\label{fig:f4}
\end{figure*}

\polarsgpu{} ran \stat{coverage.polars_gpu.baseline.supported} of
\stat{coverage.polars_gpu.baseline.total} baseline and
\stat{coverage.polars_gpu.complex.supported} of
\stat{coverage.polars_gpu.complex.total} complex query--scale cells on the GPU
(Figure~\ref{fig:f4}). The unsupported queries were
\stat{coverage.polars_gpu.unsupported_list}; the number of failures not
pre-declared was \stat{coverage.polars_gpu.unexpected_failures}. \sparkgpu{} ran
\stat{coverage.spark_gpu.baseline.supported} of
\stat{coverage.spark_gpu.baseline.total} and
\stat{coverage.spark_gpu.complex.supported} of
\stat{coverage.spark_gpu.complex.total} cells with plans meeting the GPU plan
rule. Both \polarsgpu{} gaps are expression-level: the \texttt{(?s)} inline
regex flag in \q{cx\_q13}, and \texttt{cum\_count().over(...)} in
\q{cx\_window}, whose window translation in \texttt{cudf-polars} 26.10.2
accepts only a fixed set of window functions. Resource exhaustion accounted for
\stat{coverage.polars_gpu.baseline.resource_exhaustion} baseline and
\stat{coverage.polars_gpu.complex.resource_exhaustion} complex \polarsgpu{}
cells and \stat{coverage.spark_gpu.baseline.resource_exhaustion} and
\stat{coverage.spark_gpu.complex.resource_exhaustion} \sparkgpu{} cells.
With default settings (\texttt{raise\_on\_fail=False}), the engine would
instead fall back to the CPU for these queries without failing, so a user
timing them would not be timing the GPU.

\subsection{RQ5: The version-locked CPU baseline}
\label{sec:rq5}

C1 and C4 use Polars 1.44.2 on the CPU, the version the GPU engine requires,
not the current release. In the side study (right column of
Table~\ref{tab:contrasts}), Polars~2.0.0 on the CPU was
\stat{bridge_summary.both}; at SF300 the version ratio was
\ratiotext{bridge.baseline.sf300} on the baseline suite and
\ratiotext{bridge.complex.sf300} on the complex suite, and the most extreme
query was \q{\stat{bridge_summary.complex.extreme_query}}, from
\stat{bridge_summary.complex.extreme.min} to
\stat{bridge_summary.complex.extreme.max} across scales. Each estimate rests on
\stat{bridge_summary.baseline.n_rounds} accepted rounds. The effect is not
small next to the engine contrasts: C1 is the GPU speedup over the CPU engine
of the same Polars build, not over the fastest Polars CPU release available,
and where 2.0.0 is faster, C1 and C4 overstate the GPU's advantage over the
newest CPU Polars by roughly that ratio. An exploratory, post-hoc comparison
against Polars~2.0.0 (not part of any result here) found the Polars GPU
advantage smaller but still above~1
\stat{exploratory.e1.baseline.summary.above_phrase} on the baseline
suite~\citep{ejiofor2026extended}; the companion
manuscript~\citep{ejiofor2026lockin} analyzes this version lock-in and its
consequences for GPU-speedup reporting in depth.

\subsection{Secondary analyses}
\label{sec:secondary}

The extended version~\citep{ejiofor2026extended} reports these in full; we
summarize each in one sentence. \emph{Variability}: median within-session CVs
were \stat{cv.within.polars.median}\%, \stat{cv.within.polars_gpu.median}\%,
\stat{cv.within.spark_cpu.median}\% and \stat{cv.within.spark_gpu.median}\%
(\polarscpu{}, \polarsgpu{}, \sparkcpu{}, \sparkgpu{}), and median between-round
CVs \stat{cv.between.polars.median}\%, \stat{cv.between.polars_gpu.median}\%,
\stat{cv.between.spark_cpu.median}\% and \stat{cv.between.spark_gpu.median}\%.
\emph{Order}: the largest absolute deviation of mean log-time was
\stat{order.position.max_abs_pct}\% by session position and
\stat{order.carryover.max_abs_pct}\% by predecessor engine, small next to the
engine contrasts. \emph{Startup}: \sparkgpu{}'s median launch-to-ready time
(\stat{startup.spark_gpu.median_s}\,s, outside the timer) exceeded the sum of
its six SF1 baseline query medians
(\stat{legacy.sum_medians.spark_gpu.baseline.sf1}\,s). \emph{GPU memory}: at
SF300, \polarsgpu{}'s sampled device memory (a memory-pool high-water mark)
reached \stat{peak_gpu_mem_gib.polars_gpu.complex.sf300.max}\,GiB on the
complex suite, and \sparkgpu{}'s largest task-reported allocation
\stat{task_gpu_mem_gib.spark_gpu.complex.sf300.max}\,GiB. \emph{Telemetry}: on
the complex suite at SF300, \sparkgpu{}'s samples were GPU-idle
\stat{gpu_tel.spark_gpu.complex.sf300.idle}\% of the time against
\stat{gpu_tel.polars_gpu.complex.sf300.idle}\% for \polarsgpu{}. \emph{Plan
variability}: \stat{plan_variability.cells_variable} of
\stat{plan_variability.cells_checked} Spark query--scale cells had trials that
ended in more than one final adaptive plan.

% ======================================================================
\section{The GPU-Generation Ladder}
\label{sec:ladder}
% BEGIN-RAN:any
\ifladderran
Campaigns ran on \stat{hosts.ran_list} (\stat{hosts.n_ran} hosts); hosts not
run: \stat{hosts.not_run_list}. Each campaign is analyzed on its own, exactly
as the H100 campaign above. Table~\ref{tab:ladder} sets the SF300 contrasts and
crossovers side by side, and Figure~\ref{fig:f11} compares each engine's times
with the H100. Across the \stat{ladder.hosts_count} hosts, \stat{ladder.outputs_failed_total}
of \stat{ladder.outputs_compared_total} output comparisons disagreed with
\polarscpu{}, \stat{ladder.resource_exhaustion_total} GPU-engine cells were
lost to resource exhaustion, and \stat{ladder.unexpected_failures_total}
\polarsgpu{} failures were not pre-declared.

\begin{table*}[t]
\centering
\caption{Within-host contrasts at SF300 on each host of the ladder:
suite-level geometric-mean speedup ratios with 95\% CIs ($t$, df~=~3, across
that host's four rounds). Each row comes from one host's campaign alone; rows
are not pooled. Crossover: the smallest scale at which the C4 interval lies
entirely above~1. \hostnotrun{}: campaign not run. \chartnote}
\label{tab:ladder}
\small
\setlength{\tabcolsep}{5pt}
\begin{tabular}{@{}llccccc@{}}
\toprule
Suite & Host & C1 \polarscpu{}/\polarsgpu{} & C2 \sparkcpu{}/\sparkgpu{} & C3 \sparkgpu{}/\polarsgpu{} & C4 \polarscpu{}/best GPU & Crossover \\
\midrule
baseline & A100  & \hostratioci{a100}{c1.baseline.sf300} & \hostratioci{a100}{c2.baseline.sf300} & \hostratioci{a100}{c3.baseline.sf300} & \hostratioci{a100}{c4.baseline.sf300} & \hoststat{a100}{crossover.baseline.c4.scale} \\
         & H100  & \ratioci{c1.baseline.sf300} & \ratioci{c2.baseline.sf300} & \ratioci{c3.baseline.sf300} & \ratioci{c4.baseline.sf300} & \stat{crossover.baseline.c4.scale} \\
         & GH200 & \hostratioci{gh200}{c1.baseline.sf300} & \hostratioci{gh200}{c2.baseline.sf300} & \hostratioci{gh200}{c3.baseline.sf300} & \hostratioci{gh200}{c4.baseline.sf300} & \hoststat{gh200}{crossover.baseline.c4.scale} \\
         & B200  & \hostratioci{b200}{c1.baseline.sf300} & \hostratioci{b200}{c2.baseline.sf300} & \hostratioci{b200}{c3.baseline.sf300} & \hostratioci{b200}{c4.baseline.sf300} & \hoststat{b200}{crossover.baseline.c4.scale} \\
\midrule
complex  & A100  & \hostratioci{a100}{c1.complex.sf300} & \hostratioci{a100}{c2.complex.sf300} & \hostratioci{a100}{c3.complex.sf300} & \hostratioci{a100}{c4.complex.sf300} & \hoststat{a100}{crossover.complex.c4.scale} \\
         & H100  & \ratioci{c1.complex.sf300} & \ratioci{c2.complex.sf300} & \ratioci{c3.complex.sf300} & \ratioci{c4.complex.sf300} & \stat{crossover.complex.c4.scale} \\
         & GH200 & \hostratioci{gh200}{c1.complex.sf300} & \hostratioci{gh200}{c2.complex.sf300} & \hostratioci{gh200}{c3.complex.sf300} & \hostratioci{gh200}{c4.complex.sf300} & \hoststat{gh200}{crossover.complex.c4.scale} \\
         & B200  & \hostratioci{b200}{c1.complex.sf300} & \hostratioci{b200}{c2.complex.sf300} & \hostratioci{b200}{c3.complex.sf300} & \hostratioci{b200}{c4.complex.sf300} & \hoststat{b200}{crossover.complex.c4.scale} \\
\bottomrule
\end{tabular}
\end{table*}

\begin{figure*}[t]
\centering
\studyfig{f11_gpu_ladder}{0.78\textwidth}
\caption{The GPU-generation ladder: suite geometric-mean time ratio
$t_{h}/t_{\mathrm{H100}}$ (log axis) for each engine (columns) and suite
(rows), with hosts in generation order and the H100 as the reference at~1.
Within each host, one marker per scale factor, SF1 to SF300 from left to right
(lighter to darker); bars are Welch-type 95\% intervals. Open markers: fewer
queries compared than at the other hosts. Below~1: faster than on the H100.
Descriptive only; the hosts differ in more than the GPU
(Table~\ref{tab:hosts}). \chartnote}
\Description{A grid of panels, one per engine and suite, showing time ratios
to the H100 for each host and scale on a logarithmic axis.}
\label{fig:f11}
\end{figure*}

\paragraph{The decision pattern recurred.} C4's baseline interval lay below~1
\stat{ladder.c4.baseline.below_text}, and above~1 from each host's crossover
scale onwards; on the complex suite it lay above~1
\stat{ladder.c4.complex.above_text}. \polarsgpu{} was faster than \sparkgpu{} in
\stat{ladder.c3.cells_above} of \stat{ladder.c3.cells} suite--scale cells
across the ladder (C3 from \stat{ladder.c3.min} to \stat{ladder.c3.max}).
% BEGIN-RAN:a100
\ifaonehundredran
On the A100, the four-engine result that the earlier three-engine study could
not measure shows the best GPU engine ahead of Polars CPU from
\stat{a100.crossover.c4.text}; that study's conclusion that Polars CPU kept
the lower aggregate time at every scale was specific to the engines it
compared. The 40\,GB device did not cost coverage, although \polarsgpu{}'s
sampled device memory reached
\stat{a100.peak_gpu_mem_gib.polars_gpu.complex.sf300.max}\,GiB of
\stat{a100.host.gpu_memory_gib}\,GiB on the complex suite at SF300.
\fi
% END-RAN:a100
% BEGIN-RAN:b200
\ifbtwohundredran
On the B200, C4 at SF300 was \stat{b200.c4.baseline.sf300.vs_h100} on the
baseline suite and \stat{b200.c4.complex.sf300.vs_h100} on the complex suite,
with the crossover at \stat{b200.crossover.c4.text}.
\fi
% END-RAN:b200

\paragraph{Newer GPUs did not consistently help.} Across hosts
(Figure~\ref{fig:f11}),
% BEGIN-RAN:a100
\ifaonehundredran
the GPU engines were slower on the A100 than on the H100, with
$t_{\mathrm{A100}}/t_{\mathrm{H100}}$ from \stat{ladder.xhost_gpu.a100.min} to
\stat{ladder.xhost_gpu.a100.max};
\fi
% END-RAN:a100
on the \stat{ladder.xhost_gpu.later.hosts} the GPU engines' ratios ranged from
\stat{ladder.xhost_gpu.later.min} to \stat{ladder.xhost_gpu.later.max}.
% BEGIN-RAN:b200
\ifbtwohundredran
\polarsgpu{} on the complex suite was slower on the B200 than on the H100
(interval above~1 at \stat{xhost.b200.polars_gpu.complex.summary.above}, up to
\stat{xhost.b200.polars_gpu.complex.summary.max}).
\fi
% END-RAN:b200
The CPU engines followed a different pattern:
% BEGIN-RAN:gh200
\ifghtwohundredran
\polarscpu{} was faster on the GH200 (ratios from
\stat{xhost.gh200.polars.overall.min} to \stat{xhost.gh200.polars.overall.max}),
with 26 threads on \hoststat{gh200}{host.cpus_arg_cores_text} physical cores
against \stat{host.cpus_arg_cores_text} on the H100. Because C1 and C4 use each
host's own CPU engine, the GH200's C4 at SF300 was
\stat{gh200.c4.baseline.sf300.vs_h100} on the baseline suite and
\stat{gh200.c4.complex.sf300.vs_h100} on the complex suite.
\fi
% END-RAN:gh200
Telemetry on every host shows \polarsgpu{} keeping the device busy at SF300
(GPU-idle share at most \stat{gpu_tel.ladder.polars_gpu.sf300.idle.max}\%) at
only \stat{gpu_tel.ladder.polars_gpu.sf300.power.min}--\stat{gpu_tel.ladder.polars_gpu.sf300.power.max}\%
of the board power limit, which is consistent with, but does not establish, work
that would gain little from a faster GPU. Each rung changes the whole host, so
we do not attribute these differences to the GPU generation; the companion
manuscript~\citep{ejiofor2026lockin} analyzes the core-count confound.
\fi
% END-RAN:any
% BEGIN-NOTRUN:any
\ifladderran\else
None of the additional ladder campaigns ran (reasons in
Section~\ref{sec:hosts}). As pre-registered, no software versions were changed
to make them run, so all results in this paper come from the H100 campaign.
\fi
% END-NOTRUN:any

% ======================================================================
\section{Discussion and Lessons}
\label{sec:discussion}

\paragraph{When does a GPU help?} On this workload the answer depends on scale
and query shape. At the smallest scale both GPU engines lost to their own
framework's CPU engine (C1 below~1 at \stat{c1.baseline.summary.below}, C2 at
\stat{c2.baseline.summary.below}), which is consistent with fixed per-query
costs that a small input cannot amortize. \polarsgpu{} was faster than Polars
1.44.2 on the CPU \stat{c1.both.above_text}. The largest per-query gains were
on grouped aggregations and joins and the smallest on the selective
scan-and-sum \q{q06}, matching the Polars guide's statement that grouped
aggregations and joins benefit most and I/O-bound queries
least~\citep{polars_gpu_guide}, although our design does not isolate query
shape as a factor.

\paragraph{Against a strong CPU baseline, gains are multiples of a few.} C4
asks whether the faster GPU engine beats Polars on the CPU of the same host, a
different question from the within-framework speedups C1 and C2. On the H100
it did so from \stat{crossover.baseline.c4.scale} (baseline) and
\stat{crossover.complex.c4.scale} (complex), reaching
\ratiotext{c4.baseline.sf300} and \ratiotext{c4.complex.sf300} at SF300. These
are multiples of a few, not orders of magnitude, in line with the expectation
of \citet{shanbhag2020gpucpu} that well-implemented CPU code bounds GPU
speedups near the bandwidth ratio; with Parquet reads and device-to-host
transfers inside the timer, the transfer costs analyzed by
\citet{lutz2020pump} are part of every GPU time here. The largest single C4
value at SF300, \ratiotext{c4.complex.sf300.q.cx_window} for \q{cx\_window}
(run by \sparkgpu{}, since \polarsgpu{} rejects it), is on the query where
Polars~2.0.0 gained most on the CPU; it comes largely from the slow
Polars~1.44.2 CPU path and is not a measure of GPU performance.

\paragraph{The denominator is part of the result.} The version lock (RQ5) is
itself a practical cost of choosing the GPU engine: it ties the CPU side to an
older Polars, and the newest CPU release removes part of the measured GPU
advantage. A GPU speedup should therefore be reported with the CPU version it
is measured against, and compared against the newest CPU release as
well~\citep{ejiofor2026lockin}.

\paragraph{Two routes to the GPU.} Both GPU engines execute through libcudf
(in different releases), so the C3 gap may reflect what surrounds the device
work: Spark's JVM, task scheduling, shuffle and plan translation, which are
inside our timer, against a single-process Polars plan. We did not measure
these components separately; the higher GPU-idle share of \sparkgpu{}
(Section~\ref{sec:secondary}) is compatible with, but does not establish, such
an overhead. \sparkgpu{} has the coverage advantage: it ran every complex
query, including the two that \polarsgpu{} rejects, and won most of the complex
cells \polarsgpu{} did not win.

\paragraph{Strict coverage and integrity checks are necessary.}
\stat{coverage.polars_gpu.complex.unsupported} of the twelve queries would have silently fallen back to the CPU under default Polars GPU
settings, and a hybrid Spark plan would have been timed as a GPU run without
plan inspection. Output agreement across engines, the no-fallback guard and the
plan rule cost little and make each ``GPU'' number auditable, and repeated,
order-balanced rounds are what turn single timings into intervals.

\paragraph{Practical guidance.} Within the limits of this study: for data at
or above the crossover scale on a host like these, try the Polars GPU engine
first and run it once with \texttt{raise\_on\_fail=True} to find queries it
would silently hand back to the CPU; for small inputs Polars on the CPU was
faster; compare against the newest Polars on the CPU, not only the version the
GPU engine supports; and if a query is not supported by the Polars GPU engine,
Spark with the cuDF plugin may run it on the GPU. On our ladder, a newer GPU
generation alone did not change these decisions.

% ======================================================================
\section{Threats to Validity}
\label{sec:threats}

\paragraph{Internal.} Four rounds give $\text{df}=3$ and a $t$ multiplier of
about 3.18, so intervals are wide; an interval including~1 does not show
equivalence. The Williams square balances position and first-order carryover
but does not remove them, and higher-order carryover is not balanced. Three
trials share one process, JIT and memory-pool state, so we use their median as
one observation per round. GPU and CPU clocks were recorded but not locked.
Files are preheated but page-cache residency is not verified; at SF300 the
inputs approach half of host RAM. Spark's adaptive execution can end a query in
different final plans in different trials. Taking the faster of two noisy GPU
times biases C4 towards the GPU side, most when the two engines are close.

\paragraph{Construct.} The Spark timer includes collecting results into Python
\texttt{Row} objects, while the Polars timer stops at a host DataFrame; result
sets are small, so we expect the effect to be small, but it is not measured.
Float64 normalization removes decimal-arithmetic costs. The suite geometric
mean weights every query equally; it is not a once-through suite time,
throughput or TPC metric. Under \texttt{raise\_on\_fail=True}, \polarsgpu{} does
not time a query it cannot fully run on the GPU, whereas a default user would
see a CPU fallback. The Spark plan rule checks operator selection, not the
fraction of runtime spent on the GPU. Session start-up is outside the timer.

\paragraph{External.} The primary results come from one cloud H100 instance
with \stat{host.vcpus} vCPUs and local NVMe; they do not describe distributed
Spark clusters, object storage or equal-cost CPU instances, and a CPU-only host
of the same price could have many more cores. The CPU reference in C1 and C4 is
an older Polars than the current release (Section~\ref{sec:rq5}). The Spark
settings are the earlier study's, not tuned for each host or scale, and the
Polars GPU engine uses defaults; tuning could change any contrast in either
direction. Twelve selected or custom queries on DuckDB-generated data are not
TPC-H and do not represent all analytical workloads. Results describe specific
releases as of October 2026, and the Polars GPU engine is in open beta. The two
coverage gaps make the C1 and C3 complex-suite estimates average over a
different query set than C2 and C4.

\paragraph{Ladder.} Between rungs the GPU, device memory, CPU model and
instruction set, RAM, CPU--GPU link and driver all change, so cross-host ratios
describe hosts, not GPUs. Equal engine threads are not equal physical cores
where hosts differ in simultaneous multithreading (Table~\ref{tab:hosts}), and
the 26-thread limit understates what the CPU engines could do on hosts with
more vCPUs; on the GH200 especially, a CPU engine using all vCPUs could narrow
or reverse a GPU advantage. The campaigns ran at different times on provider
hardware we do not control. Every campaign used the same author, code and
software, so the ladder tests robustness to hardware, not independent
replication.
% BEGIN-RAN:b200
\ifbtwohundredran
On the B200, the known Blackwell compiler defect (amendment A4) is guarded only
by the output-agreement gate: \stat{b200.agreement_cells.polars_gpu.within_gate}
of \stat{b200.agreement_cells.polars_gpu.total} \polarsgpu{} and
\stat{b200.agreement_cells.spark_gpu.within_gate} of
\stat{b200.agreement_cells.spark_gpu.total} \sparkgpu{} suite--scale cells had
no output outside the tolerance (largest relative differences
\stat{b200.agreement_max.polars_gpu.max_rel_diff} and
\stat{b200.agreement_max.spark_gpu.max_rel_diff}). This is absence of evidence
that the defect affected these outputs, not proof that it did not; a defect
that only slows a query would not be detected.
% NOTE: if any within_gate count is below its total, the B200 round
% failed integrity; rewrite this sentence around that outcome.
\fi
% END-RAN:b200

% ======================================================================
\section{Artifact and Extended Version}
\label{sec:artifact}

The code, saved results and analysis needed to reproduce this study are
public at \url{\artifacturl}.\artifactdoisentence{} The repository contains the generators and input verifier,
the queries (SQL and native Polars), the runner, the analyzers and plotting
code, the locked design with its amendments, and for every host the saved
results: every trial time and output, Spark final plans, Polars GPU guard
records, telemetry, manifests and checksums.
Generated Parquet data and NVIDIA binaries are not redistributed; the
generation recipe and pinned dependency versions (with the plugin JAR's
SHA-256) are provided, together with the scripts that provision a host,
generate and verify the inputs, run a campaign and rerun the analysis. Every
number in this paper is inserted from the analyzer's \texttt{stats.json}
output rather than typed. \texttt{REPRODUCE.md} gives step-by-step
instructions to rebuild the papers, re-run the analysis from the archived raw
data, and re-run the experiment.
% TODO(zenodo-doi): set \zenododoi in pvldb-macros.tex once the archive exists.

The extended version of this paper~\citep{ejiofor2026extended} contains what
is cut here for length: the full design text (including the Spark
configuration and the session lifecycle), per-query medians, legacy sums and
per-query forest plots at every scale, the variability, order, memory and
telemetry analyses, the exploratory post-hoc comparisons against Polars~2.0.0,
per-host paragraphs, and full contrast tables for every ladder host. The
companion manuscript~\citep{ejiofor2026lockin} is a secondary analysis of the
same campaigns that focuses on the version lock-in of the CPU baseline and the
core-count confounds of cross-host comparison.

% ======================================================================
\section{Conclusion}
\label{sec:conclusion}

We designed and pre-registered a four-engine, four-round randomized comparison
of CPU and GPU variants of Polars and Spark on one H100 host, with integrity
gates that do not depend on timings and interval estimates for every contrast.
On the H100, \polarsgpu{} was slower than Polars~1.44.2 on the CPU at
\stat{c1.baseline.summary.below} and faster \stat{c1.both.above_text}; the best
GPU engine was faster than Polars~1.44.2 on the CPU from
\stat{crossover.baseline.c4.scale} (baseline) and
\stat{crossover.complex.c4.scale} (complex), by \ratiotext{c4.baseline.sf300}
and \ratiotext{c4.complex.sf300} at SF300. \polarsgpu{} was the faster GPU
engine everywhere but rejected \stat{coverage.polars_gpu.complex.unsupported}
of the twelve queries as unsupported (CPU fallback disabled by design).
Polars~2.0.0 on the CPU, which the GPU engine cannot use, was
\stat{bridge_summary.both}.
% BEGIN-RAN:any
\ifladderran
Repeating the campaign on \stat{hosts.other_ran_list} did not change that
answer qualitatively: the GPU engines were slower on the A100 and not
consistently faster on the GH200 or B200 than on the H100.
\fi
% END-RAN:any
These results describe specific hosts at 26 engine threads, one derived
workload and specific software releases; they do not establish a universal
winner.

\begin{acks}
% TODO(author): funding and compute provider, if any.
No endorsement by NVIDIA, the Polars project, DuckDB, the Apache Software
Foundation or TPC is implied.
\end{acks}

\bibliographystyle{ACM-Reference-Format}
\bibliography{references}

\end{document}
